\section{Conclusion}
\label{sec: conclusion}

We addressed the question of \textit{why} stochastic tangents destabilize Mean Flow learning and provided a unifying theoretical answer: the conditional velocity field plays two distinct statistical roles in the MeanFlow loss---unbiased Monte Carlo target and Monte Carlo control variate inside a semi-gradient JVP---and the original loss assigns the wrong coefficient to the latter. We derived the optimal control-variate coefficient $\beta^{\star}$ in closed form and showed that vanilla MeanFlow's choice ($\beta\!=\!0$) is generally suboptimal because per-step gradient variance scales unboundedly with $\boldsymbol{g}^{\!\top}\!\mJ\Sigma_{\vv'}\mJ^{\!\top}\boldsymbol{g}$. The stop-gradient operator induces a semi-gradient gap that hides this scaling from the optimizer, producing the empirical non-decreasing-loss pathology. Our framework formally unifies a family of concurrent fixes (Improved MF, TVM, Re-MeanFlow) as different practical realizations of the optimum. To validate the framework empirically, we instantiated the $\beta\!=\!1$ corner via an EMA proxy with a flow-matching anchor and ran a controlled $\beta$-sweep ($\beta\!\in\!\{0,0.25,0.5,1\}$) at toy and DiT-B/4 / ImageNet-$256$ scale (\Cref{appx: dit-results}). The sweep recovers the predicted $3\!\times\!\!\sim\!5\!\times$ gradient-variance reduction with up to $-54\%$ SW$_{1}$ on the high-curvature \texttt{swiss\_roll} dataset, and recovers the predicted bias-variance FID ordering at every matched-step DiT checkpoint. The same DiT measurement also exposes a quantitative \emph{FID-vs-MSE landscape mismatch}: FID-axis bias amplification ($\sim\!1.4{:}1$ across $\beta\!=\!0.5/\beta\!=\!1$) is sublinear in the MSE-axis coefficient ($1{:}4$), so the FID-optimal $\beta$ shifts beyond the interior gradient-MSE minimizer to the unbiased corner.

\paragraph{Limitations.} The optimal coefficient~\Cref{thm: control-variate-optimum} is derived under a scalar-isotropic approximation of $\mJ$ and $\Sigma_{\vv'}$; in the full matrix-valued setting the optimum is direction-dependent and a single global $\beta$ leaves residual variance on the table. Our DiT-B/4 evaluation is single-seed; multi-seed FID confirmation will be added in the camera-ready once the runs complete. The FID-vs-MSE landscape mismatch we identify suggests that the right level at which to optimize for FID is not the gradient-MSE that Theorem~4 controls; characterizing the FID-axis loss landscape directly is left to future work.

\paragraph{Future work.} Two directions stand out. \textit{(i) Beyond MeanFlow.} The control-variate diagnosis applies to any self-supervised one-step generative model that bootstraps a parametrized map through a JVP-style total derivative, including consistency models~\cite{song2023consistencymodels} and shortcut models~\cite{frans2025stepdiffusionshortcutmodels}. We expect the same variance reduction to be available there. \textit{(ii) Adaptive $\beta$.} A per-sample coefficient driven by online estimates of $\kappa$ and $\|\vb\|$ would close the gap between our scalar optimum and the matrix-valued one, with potential gains on heterogeneous data.
